\documentclass[11pt,a4paper]{article}
\usepackage[T1]{fontenc}
\usepackage[utf8]{inputenc}
\usepackage{amsmath,amssymb,amsthm}
\usepackage[margin=1in]{geometry}
\usepackage[colorlinks=true,linkcolor=blue,urlcolor=blue]{hyperref}
\theoremstyle{plain}
\newtheorem{theorem}{Theorem}
\newtheorem{lemma}[theorem]{Lemma}
\newtheorem{proposition}[theorem]{Proposition}
\newtheorem{corollary}[theorem]{Corollary}
\theoremstyle{definition}
\newtheorem{remark}[theorem]{Remark}

\title{The empirical recurrence for OEIS A183629, proved by transfer matrix}
\author{Adrian Perez Fontelles\\ \small Independent researcher}
\date{31 August 2026}
\begin{document}
\maketitle

\begin{abstract}
OEIS A183629 counts $7$-column arrays over $\{0,\dots,2\}$ in which every
$2\times2$ subblock sums to one of $\{4\}$, and carries an empirical
recurrence of order $3$ contributed by R.~H.~Hardin. It is true, and it is not an
empirical matter at all: the rows of such an array are the vertices of a finite digraph in
which the admissible row-to-row steps are the edges, so $a(n)=\mathbf{1}^{\!\top}M^{n}
\mathbf{1}$ for the adjacency matrix $M$ of that digraph, on $S=2187$ vertices. Writing
$q$ for the characteristic polynomial of the conjectured recurrence, the recurrence holds
for every $n$ exactly when $\mathbf{1}^{\!\top}M^{j}q(M)\mathbf{1}=0$ for all
$j\ge0$, and the Cayley--Hamilton theorem bounds that to $j<S$. It is therefore a finite
computation in exact integer arithmetic, and it comes out zero.
\end{abstract}

\noindent\small 2020 Mathematics Subject Classification. 05A15, 11B37, 15B36.\normalsize

\section{The sequence and the conjecture}

OEIS A183629 is ``Number of (n+1) X 7 0..2 arrays with every 2 X 2 subblock summing to 4.''. It has offset $1$
and begins
\[
2445, 2967, 4029, 6207, 10725, 20247, 40749, 86127,\ \dots
\]
The entry carries the following comment:
\begin{quote}\small
Empirical: a(n) = 6*a(n-1) - 11*a(n-2) + 6*a(n-3).
\end{quote}
As of the ``Last modified'' line on the live entry (Mar 31 2018 14:33:31, revision 12) the statement is
still recorded as empirical, and nothing on the entry records it as settled either way.

Written out, the claim is
\begin{equation}\label{eq:conj}
a(n)\;=\;6\,a(n-1) - 11\,a(n-2) + 6\,a(n-3)
\end{equation}
for all $n>2$, which is every $n$ at which a recurrence of order $3$ can be
stated.

\section{The rows form a digraph}

Let $V=\{0,1,\dots,2\}^{7}$, the set of possible rows, so $|V|=S=2187$. For
$r,s\in V$ declare $r\to s$ an edge when, for every $j$ with $1\le j\le 7-1$,
\begin{itemize}
\item $r_j+r_{j+1}+s_j+s_{j+1}\in\{4\}$;

\end{itemize}
that is, when placing the row $s$ directly under the row $r$ satisfies the entry's
condition on every $2\times2$ subblock formed by those two rows. Let $M\in\{0,1\}^{S\times
S}$ be the adjacency matrix of this digraph and $\mathbf{1}$ the all-ones vector.

\begin{lemma}\label{lem:walk}
The arrays counted by the entry with $n+1$ rows are exactly the walks
$r^{(0)}\to r^{(1)}\to\cdots\to r^{(n)}$ of length $n$ in this digraph. Consequently
\[
a(n)\;=\;\mathbf{1}^{\!\top}M^{n}\mathbf{1} .
\]
\end{lemma}

\begin{proof}
An array with $n+1$ rows and $7$ columns is precisely a sequence
$r^{(0)},\dots,r^{(n)}$ of elements of $V$. Every $2\times2$ subblock of the array lies
in two vertically adjacent rows $r^{(i)},r^{(i+1)}$ and in two horizontally adjacent
columns $j,j+1$; so the array satisfies the entry's condition on all of its $2\times2$
subblocks if and only if each consecutive pair $\bigl(r^{(i)},r^{(i+1)}\bigr)$ satisfies
it on all column pairs, which is exactly the edge relation. Counting walks of length $n$ by
summing over all starting and ending vertices gives
$\mathbf{1}^{\!\top}M^{n}\mathbf{1}$, since
$\bigl(M^{n}\bigr)_{r,s}$ is the number of walks of length $n$ from $r$ to $s$.
\end{proof}

\section{The criterion}

Let $q(t)=t^{3}-\sum_{i}c_i\,t^{3-i}$ be the characteristic polynomial of
\eqref{eq:conj}, with the $c_i$ read off that equation.

\begin{lemma}\label{lem:crit}
For every $n\ge3$,
\[
a(n)-\sum_i c_i\,a(n-i)\;=\;\mathbf{1}^{\!\top}M^{\,n-3}\,q(M)\,\mathbf{1} .
\]
Hence \eqref{eq:conj} holds for every $n\ge3$ if and only if
$\mathbf{1}^{\!\top}M^{j}q(M)\mathbf{1}=0$ for every $j\ge0$; and it is enough to
check $j=0,1,\dots,S-1$.
\end{lemma}

\begin{proof}
By Lemma~\ref{lem:walk}, $a(n-i)=\mathbf{1}^{\!\top}M^{n-i}\mathbf{1}$, so
\[
a(n)-\sum_i c_i a(n-i)
=\mathbf{1}^{\!\top}\Bigl(M^{n}-\sum_i c_i M^{n-i}\Bigr)\mathbf{1}
=\mathbf{1}^{\!\top}M^{\,n-3}\Bigl(M^{3}-\sum_i c_i M^{3-i}\Bigr)\mathbf{1},
\]
which is the stated identity. The equivalence follows by letting $n-3=j$ run over
$\ge0$. For the last clause put $w=q(M)\mathbf{1}$. By the Cayley--Hamilton theorem
$M^{S}$ is an integer combination of $I,M,\dots,M^{S-1}$, so
$\mathbf{1}^{\!\top}M^{j}w$ for $j\ge S$ is the corresponding combination of the
earlier values; if those all vanish, so do all the rest.
\end{proof}

\section{The computation}

\begin{theorem}
The recurrence \eqref{eq:conj} holds for every $n>2$. This is the statement of
Section~1.
\end{theorem}

\begin{proof}
The vector $w=q(M)\mathbf{1}\in\mathbb{Z}^{S}$ was formed by $3$ matrix--vector
products in exact integer arithmetic, and $\mathbf{1}^{\!\top}M^{j}w$ was evaluated for
$j=0,1,\dots$ in exact integer arithmetic. Every one of those integers is $0$ from
$j=2-3+1$ onwards, and the run was continued past $S$ consecutive zeros, which by
Lemma~\ref{lem:crit} settles every larger $j$ as well. Hence the recurrence holds for
every $n>2$.
\end{proof}

\begin{remark}
The argument gives more than the recurrence: it shows the sequence is $C$-finite with
characteristic polynomial dividing that of $M$, so it has a rational generating function
whose denominator divides $\det(I-xM)$. The conjectured recurrence is one consequence; the
minimal one is a divisor of $q$.
\end{remark}

\section{Verification}

Three checks, each independent of the algebra above.

First, the digraph was built directly from the entry's stated condition and the walk counts
$\mathbf{1}^{\!\top}M^{n}\mathbf{1}$ compared against the entry's DATA at every one of
the $23$ published indices; the values agree exactly. That is what ties the digraph to
the sequence: a model that reproduced only some of the published terms would be the wrong
model, and would have been rejected.

Second, the recurrence \eqref{eq:conj} was evaluated on those published terms in exact
integer arithmetic, with no matrices involved, and vanishes wherever it applies.

Third, the annihilation test of Lemma~\ref{lem:crit} was run to the full Cayley--Hamilton
bound $j<S=2187$ rather than to a sampled prefix, in exact integer arithmetic, and returned
zero throughout.

\begin{thebibliography}{9}
\bibitem{oeis} The OEIS Foundation, \emph{The On-Line Encyclopedia of Integer
Sequences}, \url{https://oeis.org}, 2026. Sequence A183629.
\bibitem{stanley} R.~P.~Stanley, \emph{Enumerative Combinatorics}, Volume 1, 2nd ed.,
Cambridge University Press, 2011. (Transfer-matrix method, Section 4.7.)
\bibitem{fs} P.~Flajolet and R.~Sedgewick, \emph{Analytic Combinatorics}, Cambridge
University Press, 2009.
\bibitem{horn} R.~A.~Horn and C.~R.~Johnson, \emph{Matrix Analysis}, 2nd ed., Cambridge
University Press, 2013.
\end{thebibliography}

\end{document}
